تستخرج نماذج الرؤية التأسيسية مثل ``دينو في 2''~\cite{dinov2} و``كليب''~\cite{radford2021clip}
سمات دلالية قوية، لكن بدقة مكانية منخفضة جدًا، مما يحدّ من استخدامها في مهام التنبؤ الكثيف.
وتعالج مُرقّيات السمات مثل ``فيت أب''~\cite{featup} و``لوفت أب''~\cite{loftup} هذه المشكلة
باستعادة خريطة سمات عالية الدقة من أخرى منخفضة الدقة دون تشويه معنى السمات، ويقوم النموذج
المتطور ``إيني أب''~\cite{anyup} بذلك عند أي دقة خرج ومع أي مُرمِّز. غير أن ``إيني أب'' صُمِّم
للصور الثابتة؛ فعند تطبيقه على الفيديو يرقّي كل إطار على حدة ويتجاهل علاقة الإطارات ببعضها
عبر الزمن، فيصبح خرجه وميضيًا وتتزحزح حدود الأجسام من إطار إلى آخر. تقترح هذه الأطروحة
``إيني أب ثري دي''، وهو امتداد لـ``إيني أب'' يضيف الاتساق الزمني مع الحفاظ على قدرته على
العمل مع أي عمود فقري. ويُدخل تغييرين على خط أنابيب الترقية: تصبح تلافيف المُرمِّز ثلاثية
الأبعاد ليقرأ فترات زمنية قصيرة إلى جانب المكان؛ ويصبح الانتباه المتقاطع ثلاثي الأبعاد،
فيستند كل بكسل في الخرج إلى الإطارات المجاورة إضافةً إلى البكسلات المجاورة عند تحديد سمته.
دُرِّب ``إيني أب ثري دي'' على مجموعة جزئية من ``يو سي إف-101'' (1{,}170 مقطعًا، 10 أصناف
حركية)، بدءًا بالإطارات المفردة ثم بزيادة طول المقطع تدريجيًا. ويستخدم التدريب دالة خسارة
رباعية: إعادة البناء مقابل المُرمِّز المعلّم، والتوافق مع الدخل منخفض الدقة، والثبات تحت
اضطراب الصورة، وحدّ زمني يعاقب الوميض لكنه يتعطّل حين تختلف الإطارات المتتالية اختلافًا
حقيقيًا، حتى لا تُكبَح الحركة الفعلية. وعلى ``دافيس-2017''، يرتفع تشابه السمات بين الإطارات
المتجاورة من $0.863$ إلى $0.944$ ($+0.081$) مقارنةً بالخط الأساس ثنائي الأبعاد القائم على
معالجة كل إطار على حدة، مع تحسّن في المقاطع الثلاثين جميعها. وفي اختبار نشر التسميات دون
تدريب ودون رأس تجزئة متعلَّم، يرتفع مقياس ``جيه آند إف'' من $0.266$ إلى $0.347$ ويتفوّق
``إيني أب ثري دي'' في 27 مقطعًا من 30 --- رغم تطبيقه دون إعادة تدريب على سمات ``محوّل سوين
للفيديو''، وهو مُرمِّز لم يره أثناء التدريب. يسهم هذا العمل بتقديم أول مُرقٍّ للسمات متسق
زمنيًا وغير مرتبط بعمود فقري بعينه ولا بمهمة بعينها، إضافةً إلى خط أنابيب تدريب قابل
لإعادة الإنتاج. وقد تحققت هذه النتائج بـ6{,}000 خطوة تدريب فقط، بعد تهيئة الأوزان انطلاقًا
من الأوزان الرسمية لـ``إيني أب تو دي''.